\documentclass[11pt]{article}
\usepackage[margin=1in]{geometry}
\usepackage{url}
\usepackage{parskip}
\setlength{\parindent}{0pt}

\begin{document}
\thispagestyle{empty}

Dear Editors,

Please find enclosed our submission entitled

\begin{center}
\textbf{``Boyd's conductor-11 Mahler measure conjecture: proof of the
split-integral identity (C3), with an exact structural analysis of the
family $S_k$''}
\end{center}

by Huimin Zheng (College of Information and Network Engineering, Anhui
Science and Technology University, Fengyang, Anhui 233100, P.~R.~China;
\url{zhhm@ahstu.edu.cn}), which we submit for consideration in
\emph{Experimental Mathematics}.

This manuscript is original, has not been published previously, and is
not under consideration by any other journal. All authors (a single
author) are aware of and approve this submission.

\textbf{Summary and significance.} This is a story of experimental
mathematics carried through to a rigorous proof. Boyd's 1998 numerical
experiments produced conjectural identities between Mahler measures of
two-variable polynomials and $L$-values of elliptic curves, beginning
with the smallest possible conductor, $N=11$. Of the three
conductor-11 identities, the first two were proved by Brunault; the
third, (C3), concerns a polynomial vanishing on the unit torus and
asserts that a signed split integral of $\log|y|$ around the branch
cut equals $b_{11}=L'(E_{11},0)$. This paper proves (C3), completing
the conductor-11 block of Boyd's table. Highlights:
\begin{itemize}
\item \textbf{Computation-driven proof}: the proof strategy --- closing
Samart's signed open chain by a small-branch compensating arc --- was
found through numerical experimentation; the crucial homology
statement (period ratio equal to exactly $2$) is then pinned down by
certified interval arithmetic (Arb ball arithmetic with rigorous error
bounds), not merely observed numerically.
\item \textbf{Full reproducibility}: every certification step is carried
out within interval arithmetic, and all scripts, frozen certificates,
and run logs are openly available at
\url{https://github.com/huiminZheng-collab/boyd-conductor11} and
archived at \url{https://doi.org/10.5281/zenodo.21820650}.
\item \textbf{New experimental results}: a complete structural analysis
of the family $S_k$, new Boyd-type numerical identities (recorded as
conjectures, to up to $70$ digits, PARI/GP cross-checked), and a
mechanistic analysis of the conductor-53 case showing precisely why
the closed-cycle/modular-unit mechanism provably fails there --- a
negative result that delimits the method.
\item A direct regulator computation via Brunault's proved Siegel-unit
formula bypasses Bloch's diamond theorem and its normalization
ambiguities.
\end{itemize}

We believe this combination of conjecture discovery, certified
computation, structural classification, and mechanism-level
counterexample analysis is exactly the kind of work
\emph{Experimental Mathematics} aims to publish.

\textbf{Transparency.} The research was carried out by the author with
the assistance of the AI system Kimi (Moonshot AI); a formal
Declaration on the use of AI tools appears on the title page of the
article, and the author has verified all mathematical content and
takes full responsibility for it. A preprint has been submitted to
arXiv (identifier to be announced; available on request).

\textbf{Suggested referees.}
\begin{itemize}
\item Matilde N. Lal\'in (Universit\'e de Montr\'eal)
\item Fran\c{c}ois Brunault (ENS de Lyon)
\item Mathew Rogers
\item Wadim Zudilin (Radboud University)
\item Detchat Samart
\item Marie-Jos\'e Bertin
\end{itemize}
The author has no conflicts of interest with the suggested referees.

Thank you for your consideration.

\vspace{2em}
Sincerely,

Huimin Zheng\\
Anhui Science and Technology University\\
\url{zhhm@ahstu.edu.cn}

\end{document}
